\chapter{Software Requirements}
% ============================================================

\section{User Interactions / Use Cases}

The system has four actors:

\begin{itemize}
  \item \textbf{Validation Engineer (human).} Configures scenarios, executes runs,
        generates datasets, trains and evaluates models, inspects results.
  \item \textbf{vSOC Analyst (human).} Reviews attributed alerts, triages incidents.
  \item \textbf{Zone Controller (system).} Observes local TSN signals, extracts
        features, performs inference, emits alerts, participates in federated rounds.
  \item \textbf{Central Vehicle Compute (system).} Aggregates model parameters,
        distributes the global model, receives and logs alerts.
\end{itemize}

Primary use cases: UC-1 Configure and execute scenario; UC-2 Inject attack with
ground truth; UC-3 Extract feature dataset; UC-4 Train local model; UC-5 Aggregate
federated model; UC-6 Classify observation window; UC-7 Emit attributed alert;
UC-8 Review incident queue; UC-9 Evaluate and export metrics.

\section{Use Case Diagram}

\begin{figure}[htbp]
\centering
% ---- Use case diagram ----
\begin{tikzpicture}[
  actor/.style={draw, minimum width=1.6cm, minimum height=0.6cm, align=center, font=\small},
  uc/.style={draw, ellipse, minimum width=3.2cm, minimum height=0.85cm, align=center, font=\scriptsize},
  node distance=0.45cm and 2.0cm
]
\node[actor] (ve) {Validation\\Engineer};
\node[uc, right=2.2cm of ve, yshift=2.6cm] (uc1) {UC-1 Configure \&\\execute scenario};
\node[uc, below=of uc1] (uc2) {UC-2 Inject attack\\with ground truth};
\node[uc, below=of uc2] (uc3) {UC-3 Extract\\feature dataset};
\node[uc, below=of uc3] (uc4) {UC-4 Train\\local model};
\node[uc, below=of uc4] (uc5) {UC-5 Aggregate\\federated model};
\node[uc, below=of uc5] (uc6) {UC-6 Classify\\observation window};
\node[uc, below=of uc6] (uc7) {UC-7 Emit\\attributed alert};
\node[uc, below=of uc7] (uc8) {UC-8 Review\\incident queue};
\node[uc, below=of uc8] (uc9) {UC-9 Evaluate \&\\export metrics};

\node[actor, right=2.2cm of uc5, yshift=1.8cm] (zc) {Zone\\Controller};
\node[actor, right=2.2cm of uc7, yshift=-0.6cm] (cvc) {Central Vehicle\\Compute};
\node[actor, below=3.0cm of ve] (an) {vSOC\\Analyst};

\draw (ve) -- (uc1); \draw (ve) -- (uc2); \draw (ve) -- (uc3); \draw (ve) -- (uc9);
\draw (an) -- (uc8);
\draw (zc) -- (uc4); \draw (zc) -- (uc6); \draw (zc) -- (uc7);
\draw (cvc) -- (uc5); \draw (cvc) -- (uc7); \draw (cvc) -- (uc8);
\end{tikzpicture}
\caption{Use case diagram. Human actors are shown left, system actors right.}
\label{fig:usecase}
\end{figure}

\section{Detailed Use Cases}

\begin{longtable}{p{0.2\textwidth} p{0.74\textwidth}}
\toprule
\multicolumn{2}{l}{\textbf{UC-2: Inject attack with ground truth}} \\
\midrule
\endfirsthead
Actor & Validation Engineer \\
Precondition & A valid topology and benign baseline configuration exist \\
Trigger & Engineer selects an attack configuration and seed \\
Main flow &
(1) Engineer selects attack class, target zone, start time, intensity and seed.
(2) Simulation executes with injection active over the specified interval.
(3) Harness records, for every observation window, the attack class and the set of
affected zones.
(4) Raw per-node vectors and the ground-truth record are written to the results
directory. \\
Postcondition & A seeded, reproducible run with labels attributable to class, window
and zone \\
Exception & If the attack produces no measurable deviation from baseline, the run is
flagged for parameter review rather than silently retained \\
\bottomrule
\end{longtable}

\begin{longtable}{p{0.2\textwidth} p{0.74\textwidth}}
\toprule
\multicolumn{2}{l}{\textbf{UC-6: Classify observation window}} \\
\midrule
\endfirsthead
Actor & Zone Controller \\
Precondition & A trained global model is deployed locally \\
Trigger & A complete observation window of local TSN signals is available \\
Main flow &
(1) Controller computes the feature vector for the window.
(2) Vector is passed to the local copy of the global model.
(3) Model returns a class and a confidence value.
(4) If the class is not benign for $k$ consecutive windows, an alert is raised. \\
Postcondition & Either no action, or an attributed alert \\
Exception & If no valid model is available, the rule-based fallback detector operates
and the degraded state is made visible \\
\bottomrule
\end{longtable}

\section{Functional Requirements}

\subsection{Module 1: Simulation and Attack Injection}
\begin{itemize}
\item \textbf{FR-01.} The system shall instantiate a configurable multi-zone TSN
topology comprising a grandmaster, a backbone switch, $N$ zone controllers, endpoint
ECUs and a central vehicle compute node.
\item \textbf{FR-02.} The system shall synchronise all time-aware nodes using gPTP.
\item \textbf{FR-03.} The system shall execute reproducibly: a given configuration and
seed shall regenerate the same scenario structure and ground-truth labels.
\item \textbf{FR-04.} The system shall generate benign traffic and at least two
attack classes against TSN synchronisation.
\item \textbf{FR-05.} The system shall record, for every injected attack, the attack
class, the affected zone(s), and the start and end time of the attack interval.
\end{itemize}

\subsection{Module 2: Feature Extraction}
\begin{itemize}
\item \textbf{FR-06.} The system shall convert raw per-node simulation output into
fixed-size observation windows.
\item \textbf{FR-07.} The system shall compute, per window per zone, protocol-native
timing features including clock offset statistics, sibling offset residual, and
synchronisation interval residual.
\item \textbf{FR-08.} The system shall label each window with its attack class and
shall attribute the label to the affected zone(s) only.
\item \textbf{FR-09.} The system shall export the labelled dataset in a documented
schema with a stable column contract.
\item \textbf{FR-10.} Feature extraction shall be passive and shall not alter
forwarding behaviour in the simulated network.
\end{itemize}

\subsection{Module 3: Detection and Federation}
\begin{itemize}
\item \textbf{FR-11.} The system shall train a multiclass detection model on the
labelled feature dataset.
\item \textbf{FR-12.} The system shall support training in three configurations,
namely local-only per zone, centralized pooled, and federated (FedAvg), using an identical
model architecture and identical scenario splits.
\item \textbf{FR-13.} The federated coordinator shall receive model parameters only;
raw feature records shall not leave the client.
\item \textbf{FR-14.} The system shall classify an observation window into benign or a
specific attack class and shall return a confidence value.
\item \textbf{FR-15.} The system shall provide a rule-based fallback detector that
operates when no valid model or coordinator is available, and shall make the degraded
state visible.
\end{itemize}

\subsection{Module 4: Alerting and Evaluation}
\begin{itemize}
\item \textbf{FR-16.} The system shall emit alerts containing attack class, affected
zone, time window, confidence and supporting feature values.
\item \textbf{FR-17.} The system shall present alerts through an operator-facing
incident queue supporting filtering by attack class, severity and status.
\item \textbf{FR-18.} The system shall export machine-readable evaluation metrics
linked to the configuration, seed and model version that produced them.
\end{itemize}

\section{Non-Functional Requirements}

\subsection{Reliability}
\begin{itemize}
\item \textbf{NFR-01.} A documented configuration and seed shall regenerate the same
scenario structure and ground-truth labels on a clean environment.
\item \textbf{NFR-02.} Loss of the coordinator or of a valid model shall produce a
visible fallback state rather than silent cessation of monitoring.
\end{itemize}

\subsection{Usability}
\begin{itemize}
\item \textbf{NFR-03.} An operator shall be able to execute a scenario, extract
features and inspect results from documented commands without modifying source code.
\item \textbf{NFR-04.} Each alert shall carry sufficient attribution for an analyst to
identify the affected zone and time window without consulting raw simulation output.
\end{itemize}

\subsection{Performance}
\begin{itemize}
\item \textbf{NFR-05.} Local classification of one observation window shall complete
within one observation interval on the target workstation.
\item \textbf{NFR-06.} Detection shall occur within a bounded number of
synchronisation intervals from attack onset, and this latency shall be reported in
sync intervals rather than wall-clock seconds.
\end{itemize}

\subsection{Security}
\begin{itemize}
\item \textbf{NFR-07.} Raw feature records shall not be transmitted to the
coordinator; only model parameters and training metadata shall be exchanged.
\item \textbf{NFR-08.} Passive monitoring shall not increase protected-stream
99th-percentile latency by more than 5\% in matched enabled/disabled simulation runs.
\end{itemize}

% [[VERIFY]] NFR-08 is carried over from the approved proposal and is NOT yet
% measured. It is reported as planned in Section 9.5. Do not claim it as met.

\section{Acceptance Criteria}

\begin{longtable}{p{0.1\textwidth} p{0.84\textwidth}}
\toprule
\textbf{Req.} & \textbf{Objective acceptance check} \\
\midrule
\endhead
FR-03 & Two executions with identical configuration and seed produce byte-identical
ground-truth label files. \\
FR-05 & For each attack run, the recorded attack interval matches the interval over
which the timing deviation is observable in the raw trace. \\
FR-07 & For a benign run, the sibling offset residual remains within the documented
benign envelope for all zones across all seeds. \\
FR-08 & For an attack run targeting zone $z$, windows in zone $z$ carry the attack
label and windows in all other zones carry the benign label. \\
FR-10 & Protected-stream latency distributions from matched monitoring-enabled and
monitoring-disabled runs are statistically indistinguishable within the NFR-08 bound. \\
FR-12 & All three training configurations run to completion on identical splits and
produce comparable metric files. \\
FR-13 & Inspection of coordinator-bound messages confirms that no feature record
fields are present. \\
NFR-06 & Measured detection latency is lower than the latency at which accumulated
offset exceeds the configured schedule margin, across all seeds. \\
\bottomrule
\end{longtable}

\section{Requirements Traceability}

\begin{longtable}{p{0.08\textwidth} p{0.14\textwidth} p{0.2\textwidth} p{0.22\textwidth} p{0.22\textwidth}}
\caption{Requirements traceability matrix}\label{tab:traceability}\\
\toprule
\textbf{Req.} & \textbf{Module} & \textbf{Design element} & \textbf{Implementation} & \textbf{Test} \\
\midrule
\endfirsthead
\toprule
\textbf{Req.} & \textbf{Module} & \textbf{Design element} & \textbf{Implementation} & \textbf{Test} \\
\midrule
\endhead
FR-01 & M1 & Sec.~5.1, Fig.~\ref{fig:topology} & \texttt{CarazoneZonalNetwork.ned} & TC-01 \\
FR-02 & M1 & Sec.~5.2 & gPTP configuration in \texttt{omnetpp.ini} & TC-02 \\
FR-03 & M1 & Sec.~5.4 & Seeded run configurations & TC-03 \\
FR-04 & M2 & Sec.~5.8 & \texttt{DelayAttack}, \texttt{DriftAttack}, \texttt{SyncStarvation} & TC-04 \\
FR-05 & M2 & Sec.~5.7 & Ground-truth record & TC-05 \\
FR-06 & M3 & Sec.~5.8, Fig.~\ref{fig:activity} & \texttt{extract\_features.py} & TC-06 \\
FR-07 & M3 & Sec.~5.6 & \texttt{extract\_features.py} & TC-07 \\
FR-08 & M3 & Sec.~5.7 & \texttt{extract\_features.py} & TC-08 (DEF-01) \\
FR-09 & M3 & Table~\ref{tab:schema} & \texttt{features.csv} & TC-09 \\
FR-11--FR-15 & M4, M5 & Sec.~5.1--5.3 & Planned, Iteration 03 & Planned \\
FR-16--FR-18 & M6 & Sec.~5.1 & Planned, Iteration 04 & Planned \\
\bottomrule
\end{longtable}

% ============================================================
